This study asked whether ADHD and autism---two diagnostic concepts at the centre of contemporary debates about the popularisation of mental health language---have undergone the semantic loosening that concept creep theory predicts, and it did so in general web discourse rather than in the curated corpora on which prior evidence largely rests. The answer, on the measures used here, is that they have not. Across a frame-aware adaptation of the SIBling framework \citep{baesMultidimensionalFrameworkEvaluating2024}, the dominant empirical signature was not milder, broader, or diluted meaning, but a reorganisation of the discourse surrounding two comparatively stable concepts. This chapter first summarises the findings in relation to the research questions (Section~\ref{sec:discussion-summary}). It then develops three interpretive claims: that the results indicate semantic consolidation rather than concept creep (Section~\ref{sec:discussion-semantic-consolidation}); that the clearest change in ADHD and autism discourse over the period was compositional---a shift from disorder talk toward lived experience, accompanied by a recent positification (a positive valence shift) of clinical contexts (Section~\ref{sec:discussion-lived-experience}); and that the robustness checks expose a measurement sensitivity with implications for the wider literature on lexical semantic change (Section~\ref{sec:discussion-measurement-sensitivity}). Limitations (Section~\ref{sec:discussion-limitations}) and directions for future research (Section~\ref{sec:discussion-future-research}) close the chapter.

\section{Summary of Findings}
\label{sec:discussion-summary}

With respect to RQ1, ADHD and autism did not become jointly more prominent in sampled general web discourse between 2014 and 2026. Autism salience declined by a fitted 28.9\%, while the fitted 27.0\% increase for ADHD did not differ from zero at the uncorrected \(p < .05\) level. The comparators moved in different directions---loneliness rose, sadness fell, and frustration showed no reliable trend---so the autism decline cannot be attributed to a uniform contraction of negative-affect language in the crawl.

With respect to RQ2, the balance of framing shifted toward lived experience. Among contexts assigned to a single clear frame, the lived-experience share of ADHD discourse rose by 1.07 percentage points per year, a fitted increase from 17.5\% to 30.3\%; the corresponding autism trend was directionally similar but roughly half the size and was not retained by the AR(1) sensitivity model. The full predicted composition showed the same pattern, with clinical-only shares falling for both targets.

With respect to RQ3, neither of the hypothesised forms of concept creep was supported (Table~\ref{tab:lsc-regression-target-frames}). Contrary to the vertical-creep hypothesis, intensity did not decline in any target stratum; the only slope that differed from zero was a nominal increase in autism clinical contexts. Contrary to the horizontal-creep hypothesis, breadth under the target-aware operationalisation decreased for ADHD overall and in ADHD clinical contexts---the strongest semantic trends observed in the study---decreased in autism lived-experience contexts, and was flat for autism overall, with the nominally positive autism clinical slope not surviving the AR(1) check. Sentiment showed no linear trend in any target stratum; instead, the autocorrelation-flagged series traced U-shaped trajectories with minima between mid-2018 and mid-2019, and the AR(1) sensitivity model indicated a positive valence trend in autism clinical contexts.

With respect to RQ4, the thematic content of both discourses changed remarkably little. The 28 neighbours retained by the stability filter were diagnostic and child- or family-centred throughout: \emph{child}, \emph{disorder}, \emph{symptom}, \emph{diagnose}, and \emph{condition} anchored ADHD discourse, while \emph{child} remained the highest-similarity neighbour of autism overall and in clinical contexts, with \emph{people}, \emph{support}, \emph{life}, and \emph{family} characterising lived-experience contexts.

\section[Semantic Consolidation Rather than Concept Creep]{Semantic Consolidation Rather than Concept\texorpdfstring{\\}{ }Creep}
\label{sec:discussion-semantic-consolidation}

The central substantive finding is a double null with a directional surprise. ADHD and autism did not drift into less emotionally intense contexts, and they did not disperse across a broader range of lexical contexts; where reliable movement occurred, it ran in the opposite direction, most clearly in the sustained contraction of ADHD breadth, under the target-aware operationalisation, from .138 to .117. This pattern is unlikely to reflect measures that were simply too blunt to register change of any kind: the same instruments detected reliable trends among the comparators, including a broadening of frustration, an intensification of sadness contexts, and robust positification of frustration and loneliness. Within this corpus and this measurement scheme, the affective and contextual profiles of ADHD and autism were among the more stable series observed.

This stability is intelligible in light of \citeauthor{iacobComputationalAnalysisEmergence2026}'s \citeyear{iacobComputationalAnalysisEmergence2026} finding that long-established, codified psychological terms such as \emph{OCD} and \emph{bipolar} showed little year-to-year semantic movement on Reddit, whereas newer vernacular terms such as \emph{gaslighting} shifted substantially. ADHD and autism are institutionally anchored concepts: their reference is stabilised by diagnostic manuals, clinical gatekeeping, service categories, and educational law in ways that ordinary-language harm concepts such as \emph{bullying} or \emph{harassment} \citep{vylomovaSemanticChangesHarmrelated2021} are not. Concept creep in such categories may therefore operate primarily at the level of diagnostic criteria and category application---who is counted as an instance---rather than at the level of the contexts in which the label is used. This reading is consistent with \citeauthor{fabianoDiagnosticInflationDSM2020}'s \citeyear{fabianoDiagnosticInflationDSM2020} meta-analytic evidence that ADHD's diagnostic criteria inflated substantially across successive DSM revisions, and it suggests a distinction the concept creep literature has tended to elide: a category can expand extensionally, absorbing many more people, while the typical linguistic contexts of its name---which is all that collocate- and embedding-based measures index---remain stable or even converge. On this account, the rising prevalence documented in the epidemiological literature
\citep{zeidanGlobalPrevalenceAutism2022,mckechnieAttentiondeficitHyperactivityDisorder2023,cuiDSM5ChangesCOVID192026} need not leave the semantic trace that a naive application of concept creep theory would predict.

The ADHD breadth contraction invites a further interpretation: consolidation of a discourse genre. As ADHD moved from a specialist topic to a mainstream one, its web coverage appears to have converged on a recognisable register---symptom lists, awareness explainers, parenting and educational advice---rather than diversifying. The post-hoc diagnostics support this reading: the highest-distance ADHD contexts remained saturated with diagnostic and educational vocabulary (\emph{disorder}, \emph{impulsivity}, \emph{inattention}, \emph{learning}, \emph{dyslexia}), and the stable neighbours were uniformly diagnostic and child-related. Alternatively, the contraction may reflect the composition of the sample rather than the concept: if health articles written to a common template (e.g., for search-engine optimisation) account for a growing share of retained ADHD documents, their near-identical contexts would pull measured breadth down even if the meaning of ADHD itself were unchanged. The two possibilities are compatible and cannot be separated here. Either way, the frustration comparator's concurrent broadening indicates that the contraction is target-specific rather than a corpus-wide artefact.

The findings also bear on the evidentiary question that motivated the study. Prior evidence for creep in mental-health-related concepts derives largely from psychology abstracts and curated general-language corpora \citep{vylomovaSemanticChangesHarmrelated2021,baesSemanticShiftsMental2023,baesSemanticInflationTrauma2023}, and \citet{pislRevisitingSemanticSeverity2025} have shown that apparent semantic trends in such corpora can be artefacts of shifting discourse composition. The present study addressed both concerns---moving to general web discourse and stratifying by frame---and found no concept creep in ADHD and autism. This does not overturn earlier findings for other concepts, which concerned different terms with different institutional standing, but it does reinforce the conclusion that creep results are corpus- and composition-dependent, and that inferences about cultural change drawn from any single curated corpus should be made cautiously. At the same time, Common Crawl imposes structural constraints of its own \citep{baackCriticalAnalysisLargest2024}, a point developed in Section~\ref{sec:discussion-limitations}; the discrepancy between corpora identifies corpus dependence rather than settling which corpus better reflects the underlying culture.

The salience results sharpen this point. Autism's declining prominence in the sampled crawl sits awkwardly beside steeply rising diagnosed prevalence \citep{zeidanGlobalPrevalenceAutism2022,abs2024} and more than four million TikTok videos hashtagged by June 2026. The most plausible reconciliation is that this discourse growth is concentrated on social platforms and behind robots.txt restrictions, both largely outside Common Crawl's reach \citep{baackCriticalAnalysisLargest2024}: ``web discourse'' is not a single population, and the discourse boom visible on TikTok is not automatically visible---or even present---in web-scale archives.

\section{From Disorder Talk to Lived Experience}
\label{sec:discussion-lived-experience}

If the meanings of ADHD and autism held broadly steady, the way they are talked about did not. The most reliable change detected in the entire study concerns who and what these terms are used to discuss: ADHD discourse shifted markedly from clinical toward lived-experience framing, and autism discourse moved in the same direction from a substantially higher starting point (a fitted 35.9\% lived-experience share in 2014, against 17.5\% for ADHD). This asymmetry is theoretically coherent. Identity-first and community framings of autism were established well before the study window, whereas ADHD's popular turn---adult self-recognition, online community formation, and identity-oriented content---is more recent, consistent with qualitative evidence on ADHD online communities \citep{ginappExperiencesAdultsADHD2023}, clinical reports of self-diagnosed presentations \citep{hartnettSocialMediaADHD2024}, the sharp growth in ADHD media coverage \citep{martinChangingPrevalenceADHD2025}, and accounts of ``platformed diagnosis'' spilling outward from social media \citep{alperTikTokAlgorithmicallyMediated2025}. On this reading, ADHD is undergoing in the 2020s a reframing that autism discourse partially completed earlier, leaving autism's trend flatter and, after accounting for serial dependence, statistically unresolved.

The compositional shift also vindicates the study's frame-aware design on its own terms. Frame-specific trajectories repeatedly diverged in sign: autism clinical breadth trended nominally upward while lived-experience breadth declined; autism clinical sentiment was positive under the AR(1) model while the lived-experience slope pointed downward. An unstratified analysis would have blended these movements with the changing frame mixture, precisely the conflation that \citet{pislRevisitingSemanticSeverity2025} identified in newspaper text. The present results extend that demonstration to web-scale data and to a different kind of composition problem. For anxiety and depression, stratification must first separate clinical constructs from the ordinary states and non-clinical phenomena (e.g., everyday anxiety, or an economic depression) the same words denote \citep{xiaoHaveConceptsAnxiety2023}; ADHD and autism are diagnostic concepts by definition, so the composition that shifts is the framing of the diagnosis itself---as disorder construct or as identity and lived experience.

The sentiment results add a temporal texture that linear models missed. The flagged ADHD and autism clinical series traced U-shapes with minima between Q3 2018 and Q2 2019, implying declining valence through the mid-2010s followed by recovery through 2026, with the autism clinical positification robust to first-order autocorrelation. The turning point falls within the period in which neurodiversity entered mainstream media vocabulary, with US news coverage of neurodiversity and neurodivergent individuals increasing between 2016 and 2022 \citep{huang-horowitzExploringMediaDiscourse2025}, and in which first-person ADHD and autism content proliferated on social platforms \citep{hartnettSocialMediaADHD2024,alperTikTokAlgorithmicallyMediated2025}. The post-hoc contributors indicate that the recovery is carried by child-, family-, and support-related vocabulary against a persistent background of diagnostic and comorbidity terms. Even so, this is at most gentle destigmatisation, and it is not uniform: the appearance of \emph{dangerous} and \emph{obsession} among the late-period arousal-raising collocates in autism lived-experience contexts shows that alarmed or stigmatising discourse persists.

A further pattern deserves note. From 2018 onward, \emph{adhd} entered the leading negative-valence contributors in autism contexts, mirroring the persistent presence of \emph{autism} in ADHD contexts, and 11,795 corpus documents contained both targets. This is web-general corroboration of the ADHD--autism semantic convergence that \citet{kangConvergingRepresentationsAttentionDeficit2025} reported on Reddit from 2019: the two concepts are increasingly discussed together, as co-occurring conditions and as jointly constitutive of neurodivergent identity. The convergence also complicates target-specific measurement, since each concept increasingly forms part of the other's context.

Taken together, these results speak to the therapy-speak debate with which this project opened. The trivialisation concern holds that popularisation erodes the meaning of clinical terms, dispersing them across ever milder and more heterogeneous uses and thereby depriving diagnosed people of precise language for their experiences \citep{isern-masUnmaskingTherapyspeak2025,spencerIsntEveryoneLittle2021}. For ADHD and autism in general web discourse, the semantic preconditions of that concern are not in evidence: contexts did not become milder, did not broaden, and remained thematically anchored to diagnosis, childhood, and family. A different hermeneutical concern, however, does find support: despite the steep rise in adult diagnoses \citep{mckechnieAttentiondeficitHyperactivityDisorder2023}, the stable thematic core of both discourses remained resolutely child-centred, suggesting that late-diagnosed adults continued to encounter a web discourse organised around children's presentations rather than their own.

\section[Measurement Sensitivity and Its Implications for LSC Research]{Measurement Sensitivity and Its Implications for\texorpdfstring{\\}{ }LSC Research}
\label{sec:discussion-measurement-sensitivity}

The robustness checks carry an implication that extends beyond this study. Substituting the original SIBling resources for the refined ones changed some substantive conclusions. Under Warriner norms, ADHD intensity declined reliably, a result that, taken alone, would have been reported as vertical concept creep---whereas the NRC--VAD estimate was flat. Under MPNet sentence embeddings, the pronounced XL-LEXEME decline in ADHD breadth attenuated to a null, while autism's flat XL-LEXEME trajectory became a reliable MPNet decline. Two conclusions were common to both operationalisations and can be stated with corresponding confidence: neither target broadened horizontally, and neither showed a reliable linear sentiment trend. The vertical-creep conclusion, by contrast, is measurement-conditional and should be reported as such.

There are principled grounds for preferring the refined measures \textit{a priori}: NRC--VAD offers roughly threefold lexical coverage, multi-word entries, and higher reported reliability \citep{mohammadNRCVADLexicon2025}, and XL-LEXEME is a word-in-context model built for LSC detection that represents the target use itself rather than the surrounding sentence topic \citep{cassottiXLLEXEMEWiCPretrained2023}, a distinction that matters when the analytic object is a term's usage rather than a document's subject matter. But the \textit{a priori} argument does not dissolve the empirical problem: the SIBling dimensions, as currently operationalised, are not measurement-invariant, a possibility its authors anticipated in describing the operationalisations as first implementations \citep{baesMultidimensionalFrameworkEvaluating2024}. This observation may also help explain the mixed findings accumulating in the literature---increases in semantic intensity where creep predicts decreases \citep{xiaoHaveConceptsAnxiety2023}, reversals across corpora \citep{iacobComputationalAnalysisEmergence2026}, and composition-dependent trends \citep{pislRevisitingSemanticSeverity2025}---some portion of which plausibly reflects instrumentation rather than substance. Until the framework's dimensions are validated against human judgements of intensity, breadth, and connotation, LSC studies in this tradition would do well to treat multi-resource robustness checks not as optional supplements but as part of the primary analysis, as practised here.

\section{Limitations}
\label{sec:discussion-limitations}

The most consequential limitations concern the corpus. Common Crawl is neither a complete copy of the web nor a representative sample of it: domain selection follows harmonic centrality, adopted in 2017 (prior crawls relied mainly on externally donated URL seed lists, which had dwindled over time); major social platforms are absent; and rights-holder blocking via robots.txt has grown over precisely the years studied \citep{baackCriticalAnalysisLargest2024}. Temporal comparability is therefore imperfect, and trends---salience above all---partly reflect changes in what the crawl could see. Findings generalise to quality-filtered, English-language open-web discourse as archived by Common Crawl, not to public discourse at large, and emphatically not to the social platforms where much contemporary ADHD and autism discourse occurs.

The frame-aware design inherits classifier error. The lived-experience head performed adequately but weakest (F1 = .760, precision = .704), so lived-experience strata contain clinical admixture, which will have attenuated frame contrasts and blurred frame-specific trends. The more damaging possibility is non-stationary error: stylistic drift inflating late-period lived-experience false positives would manufacture exactly the trend observed. A period-stratified audit found no evidence of this---lived-experience precision and calibration were best in the late period, and human-authoritative labels on the adjudicated passages reproduce the rise independently of the classifier---though with only 11--23 lived-experience validation positives per period, modest error drift cannot be excluded, and the fitted ADHD rise may be somewhat smaller than the apparent 12.8 percentage points. The human gold standard also rests on a single coder: the pilot, the adjudication of flagged cases, and the 200-passage validation set reflect one annotator's judgement, so no inter-annotator agreement can be reported and validation performance is measured against an individual rather than a consensus standard. And although the ACT workflow retained human authority over labels, LLM-assisted annotation may import model-specific biases that human review of high-risk cases only partially corrects. The frame codebook, finally, is specific to ADHD and autism, so the baselines could not be frame-stratified, limiting target–baseline comparisons to unframed series.

Three further limitations should be explicit. The affective measures score lemmatised collocates with context-free lexicon values within the \(\pm\)5-token windows adopted from the SIBling operationalisation \citep{baesMultidimensionalFrameworkEvaluating2024}, and are therefore blind to negation, irony, and syntactic scope. The comparators are baselines for general negative-affect language, not matched semantic controls. And the design as a whole is observational: it can characterise how discourse changed, but not why. The inferential bounds set out in Section~\ref{sec:methods-statistical-analysis}---34 uncorrected trend tests and the minimum detectable effects implied by 13 annual observations---apply to all null results reported here.

\section{Directions for Future Research}
\label{sec:discussion-future-research}

Three extensions follow directly. First, cross-platform designs are needed to test the suggestion advanced in Section~\ref{sec:discussion-semantic-consolidation} that ADHD and autism discourse has migrated toward social platforms outside Common Crawl's reach: applying the same frame-aware measures to web, Reddit, and short-video corpora over a common window would show whether the semantic stability observed here coexists with creep in platform discourse \citep[cf.][]{devriesExploringConceptCreep2025,aragon-guevaraReachAccuracyInformation2025}. Second, the measurement sensitivity identified in Section~\ref{sec:discussion-measurement-sensitivity} motivates validation studies in which lexicon- and embedding-based intensity, breadth, and sentiment estimates are benchmarked against human judgements of the same contexts, so that resource choice can rest on demonstrated convergent validity rather than plausibility arguments. Third, the substantive account developed here generates testable expectations: if institutional codification protects diagnostic labels from creep, then frame-aware analyses of less codified neighbouring vocabulary---\emph{neurodivergent}, \emph{masking}, \emph{stimming}, \emph{executive dysfunction}---should reveal greater semantic movement than ADHD and autism themselves; and linking annual discourse measures to diagnostic and prescribing statistics would allow the temporal ordering assumed by the prevalence inflation hypothesis \citep{foulkesAreMentalHealth2023} to be examined directly. Finer-than-annual time resolution, non-English web discourse, and qualitative analysis of the contentious late-period lived-experience contexts flagged in Section~\ref{sec:discussion-lived-experience} would each add further resolution to the picture drawn here.
